\subsection{DET-18: rational endpoint weights and an ordinary determinant}
The finite character average at the examined endpoint admits an exact
ordinary Fredholm determinant after an explicit integer multiplication.
Before that multiplication, its logarithmic derivative is a scaled trace.
The distinction is a question about zero multiplicities, including the
ramification at the central point.

\paragraph{The actual finite family.}
Fix a source level $\mathcal L$ and write $o=\mathcal L.o$ and
$d_{\mathcal L}=3o+1$. Its nontrivial index set is
$I=\mathcal L.\mathrm{Nt}$, its primitive character at $i$ is $\chi_i$,
and its weight is
\[
 \alpha_i\in\{1,d_{\mathcal L}^{-1}\}.
\]
These are the definitions in
\path{Formalism/SharpRestriction/RamificationFrobenius/TotientGaloisClasses/ExplicitFormulaErrorOfHostConditions.lean},
lines 8--9 and 32--38.
The primitive-character requirement is part of \texttt{PChar}, in
\path{PrimePairSizeCondition.lean} in the same directory, lines 7--11;
the inverse remains primitive, lines 21--24.
The actual passage to the analytic finite family, with the zeta character
having coefficient one, is
\path{Development/AutomorphicArithmetic/PrimitiveCosets/LevelFamily.lean},
lines 9--25 and 27--52.
Thus the following construction applies to this actual primitive family.
It is not a statement about every possible \texttt{CharFamily}, whose
definition permits nonprimitive characters and arbitrary positive real
weights.

\paragraph{Completions and normalization.}
For a nontrivial primitive character of conductor $q$, put
\[
 \kappa_\chi=\frac{1-\chi(-1)}2,\qquad
 F_\chi(s)=
 (q/\pi)^{(s+\kappa_\chi)/2}
 \Gamma((s+\kappa_\chi)/2)L(s,\chi),
 \qquad \chi^\vee=\chi^{-1}.
\]
For the modulus-one character put
\[
 F_\zeta(s)=s(s-1)\pi^{-s/2}\Gamma(s/2)\zeta(s)=2\xi(s).
\]
These are precisely the Dirichlet completion conventions of
\fid{MAS-11}. The classical completed functional equation and
entire-order-one assertion are used as background:
\cite[(25.15.5)--(25.15.6)]{dlmf2515},
\cite[\S6.1; (6.2.2)]{kedlayaFE}, and
\cite[\S1, p.~282]{rajan1998}.
In particular,
\[
 F_\chi(s)=\epsilon_\chi F_{\chi^\vee}(1-s),
 \qquad \epsilon_\chi\epsilon_{\chi^\vee}=1,\qquad
 F_\zeta(s)=F_\zeta(1-s).
\]
The last root-number identity follows by applying the functional equation
twice to the nonzero function $F_\chi$.
The nonprincipal completions are entire; the two simple poles of the
principal completion are cancelled by $s(s-1)$, without adding zeros.
Nonvanishing at $s=1$ and the functional equation give
$F_\chi(0)F_\chi(1)\ne0$; also
$F_\zeta(0)=F_\zeta(1)=1$.
For the source's nonvanishing input see
\path{Development/AutomorphicArithmetic/DedekindLFunctions/LandauTheorems/CharFamilyLFunctionNeZeroOfOneLeRe.lean},
lines 9--20. Together with the functional equation this places the
completed zeros in $0<\Re\rho<1$.

Choose any positive integer $M$ such that
$b_i=M\alpha_i/2$ is an integer for every $i$.
The uniform choice
\[
 M=2d_{\mathcal L},\qquad
 b_i\in\{d_{\mathcal L},1\}
\]
always works. It is a sufficient choice, not a claim of minimality.
Define the unambiguously entire function
\begin{equation}\label{eq:DET18completedproduct}
 \mathscr P_M(s)=
 \left(\frac{F_\zeta(s)}{F_\zeta(0)}\right)^M
 \prod_{i\in I}
 \left(
 \frac{F_{\chi_i}(s)F_{\chi_i^\vee}(s)}
 {F_{\chi_i}(0)F_{\chi_i^\vee}(0)}
 \right)^{b_i}.
\end{equation}
No fractional powers occur in this definition. Define the meromorphic
weighted logarithmic derivative
\begin{equation}\label{eq:DET18weightedlog}
 \mathscr L_{\mathcal L}(s)=
 \frac{F_\zeta'(s)}{F_\zeta(s)}
 +\sum_{i\in I}\frac{\alpha_i}{2}
 \left(
 \frac{F_{\chi_i}'(s)}{F_{\chi_i}(s)}
 +\frac{F_{\chi_i^\vee}'(s)}{F_{\chi_i^\vee}(s)}
 \right).
\end{equation}
If $m_\chi(\rho)=\operatorname{ord}_\rho F_\chi$, with value zero
when $F_\chi(\rho)\ne0$, put
\begin{equation}\label{eq:DET18weightedmultiplicity}
 S_\rho=m_\zeta(\rho)+
 \sum_{i\in I}\frac{\alpha_i}{2}
 \bigl(m_{\chi_i}(\rho)+m_{\chi_i^\vee}(\rho)\bigr).
\end{equation}
At every zero in this union, $S_\rho>0$ and
$\operatorname{ord}_\rho\mathscr P_M=MS_\rho$.
Coincident zeros of different characters contribute additively.

\begin{proposition}
There is a diagonal normal operator $C_{\mathcal L,M}\in\mathcal S_1$
such that, for every $s\in\C$,
\begin{equation}\label{eq:DET18det}
 \det(I-s(1-s)C_{\mathcal L,M})=\mathscr P_M(s).
\end{equation}
For each unordered noncentral pair $\{\rho,1-\rho\}$ its diagonal
contains
$[\rho(1-\rho)]^{-1}$ with multiplicity $MS_\rho$.
At the central point it contains $4$ with multiplicity $MS_{1/2}/2$,
understood as zero if there is no central zero.
These multiplicities are nonnegative integers.
Away from the completed zero set,
\begin{equation}\label{eq:DET18resolvent}
 \mathscr L_{\mathcal L}(s)
 =\frac1M\frac{\mathscr P_M'(s)}{\mathscr P_M(s)}
 =\frac{2s-1}{M}
 \Tr\!\left(C_{\mathcal L,M}
 (I-s(1-s)C_{\mathcal L,M})^{-1}\right).
\end{equation}
The determinant and trace identities require no GRH. Under GRH for the
retained family, $0\le C_{\mathcal L,M}\le4I$.
\end{proposition}

\begin{proof}
Every paired factor in \eqref{eq:DET18completedproduct} is invariant
under $s\mapsto1-s$: its two functional-equation constants have product
one. Thus $\mathscr P_M(s)=\mathscr P_M(1-s)$,
$\mathscr P_M(0)=\mathscr P_M(1)=1$, and $\mathscr P_M$ is entire of
order at most one. The functional equation also gives
$m_{\chi}(1/2)=m_{\chi^\vee}(1/2)$. Each paired factor consequently
has even central order. The zeta factor has even central order by its
own symmetry. This proves $MS_{1/2}\in2\mathbb Z_{\ge0}$ without a
central nonvanishing assumption.

Set $E(z)=\mathscr P_M(1/2+z)$. Its entire power series has only even
powers. Consequently there is a unique entire function $\mathscr D_M$
with
\[
 E(z)=\mathscr D_M(1/4-z^2),\qquad \mathscr D_M(0)=1.
\]
This definition involves no choice of a square-root branch. If
$\mathfrak M_E(r)=\max_{|z|\le r}|E(z)|$, then
\[
 \max_{|v|\le r}|\mathscr D_M(v)|
 \le \mathfrak M_E\bigl(\sqrt{r+1/4}\bigr).
\]
Hence $\mathscr D_M$ has order at most $1/2$.

The map $s\mapsto v=s(1-s)$ is locally invertible at
$\rho\ne1/2$ and has local expression $v=1/4-(s-1/2)^2$ at the
central point. It follows that the distinct zeros of $\mathscr D_M$
are $v_\rho=\rho(1-\rho)$, with multiplicity $MS_\rho$ for each
unordered noncentral pair and multiplicity $MS_{1/2}/2$ at $v=1/4$.
There is no zero at $v=0$. Write the resulting multiplicity at a
distinct $v$ as $q_v$.

For completeness, the order bound gives the summability needed for an
ordinary determinant. Choose $1/2<\theta<1$. For large $R$ the
logarithm of the maximum modulus of $\mathscr D_M$ on $|v|\le2R$
is at most $C_\theta R^\theta$. Jensen's formula, using
$\mathscr D_M(0)=1$, implies
$\sum_{|v|\le R}q_v=O_\theta(R^\theta)$.
Dyadic summation now gives
\[
 \sum_{\mathscr D_M(v)=0}\frac{q_v}{|v|}<\infty.
\]
Hadamard factorization for an entire function of order less than one
has genus zero and no nonconstant exponential factor. Its constant
is fixed by $\mathscr D_M(0)=1$, so
\[
 \mathscr D_M(v)=\prod_{\mathscr D_M(a)=0}
 (1-v/a)^{q_a}.
\]
On a Hilbert space with one orthonormal coordinate for each repeated
zero $a$, let $C_{\mathcal L,M}$ be diagonal with entry $a^{-1}$.
The displayed summability proves membership in $\mathcal S_1$.
Its ordinary Fredholm product is the displayed product for
$\mathscr D_M$, proving \eqref{eq:DET18det}. This argument does not
infer trace class from eigenvalue summability for an unspecified
non-normal operator; the operator has been explicitly defined as
diagonal.

On a compact set disjoint from the zero set of $\mathscr D_M$,
the series
\[
 \sum_a \frac{q_a/a}{1-v/a}
\]
converges absolutely and uniformly. The tail denominators are
uniformly bounded away from zero since $a^{-1}\to0$; the remaining
denominators form a finite nonvanishing set. The same reasoning
gives a bounded diagonal inverse of $I-vC_{\mathcal L,M}$.
Thus termwise logarithmic differentiation gives
\[
 \frac{\mathscr D_M'(v)}{\mathscr D_M(v)}
 =-\Tr\!\left(C_{\mathcal L,M}(I-vC_{\mathcal L,M})^{-1}\right).
\]
The chain rule, $v'(s)=1-2s$, and the integer exponents in
\eqref{eq:DET18completedproduct} prove
\eqref{eq:DET18resolvent}.

If $S_{1/2}=0$, both sides of \eqref{eq:DET18resolvent} vanish at
$s=1/2$. If $S_{1/2}>0$, put $z=s-1/2$ and
$q=MS_{1/2}/2$. The central block contributes $q/z^2$ to the
resolvent trace and therefore $S_{1/2}/z$ to the right side of
\eqref{eq:DET18resolvent}. This is exactly the logarithmic
derivative's central pole. No division by $2s-1$, and no removal
of central zeros, is used.

Under GRH, $\rho=1/2+it$ and
$[\rho(1-\rho)]^{-1}=(1/4+t^2)^{-1}\in(0,4]$.
This proves the final assertion.
\end{proof}

\paragraph{The precise positivity condition.}
In fact, for this diagonal realization, GRH for the retained family
is equivalent to $0\le C_{\mathcal L,M}\le4I$.
To see the converse, write $\rho=\beta+it$ with $0<\beta<1$.
Then
\[
 \rho(1-\rho)=\beta(1-\beta)+t^2+i\,t(1-2\beta).
\]
Its real part is positive. If its reciprocal is real and positive,
either $\beta=1/2$ or $t=0$. In the latter case
$(\beta(1-\beta))^{-1}\ge4$, with equality only at $\beta=1/2$.
The upper bound $C_{\mathcal L,M}\le4I$ therefore forces the
critical line also in this case. Positivity alone would permit a
real off-center zero with eigenvalue greater than four; no such zero
is asserted to exist here.

\paragraph{The actual prime-sum transform, including all corrections.}
Put $\psi=\Gamma'/\Gamma$ and define
\[
 a_\zeta(s)=\frac1s+\frac1{s-1}
 -\frac12\log\pi+\frac12\psi(s/2),\qquad
 a_\chi(s)=\frac12\log(q/\pi)
 +\frac12\psi((s+\kappa_\chi)/2).
\]
The dual character has the same conductor and parity, so
$a_{\chi^\vee}=a_\chi$. Set
\[
 A_{\mathcal L}(s)=a_\zeta(s)+\sum_{i\in I}\alpha_i a_{\chi_i}(s),
 \qquad
 b_{\mathcal L}=1+\gamma_E+
 \sum_{i\in I}\frac{\alpha_i}{2}
 \left(\frac{L'(1,\chi_i)}{L(1,\chi_i)}
 +\frac{L'(1,\chi_i^\vee)}{L(1,\chi_i^\vee)}\right).
\]
The constant $b_{\mathcal L}$ is real by conjugation.
For exactly the scalar $G_{\mathcal L}$ of \fid{DET-09}, its
absolutely convergent transform on $\Re s>1$ satisfies
\begin{align}\label{eq:DET18actualPhi}
 \Phi_{\mathcal L}(s)
 &:=\int_0^\infty G_{\mathcal L}(e^u)e^{-su}\,du \notag\\
 &=\frac{A_{\mathcal L}(s)-\mathscr L_{\mathcal L}(s)}
 {s(s-1)}
 -\frac1{(s-1)^2}+\frac{b_{\mathcal L}}{s-1}\notag\\
 &=\frac{A_{\mathcal L}(s)-
 \frac{2s-1}{M}\Tr\!\left(C_{\mathcal L,M}
 (I-s(1-s)C_{\mathcal L,M})^{-1}\right)}
 {s(s-1)}
 -\frac1{(s-1)^2}+\frac{b_{\mathcal L}}{s-1}.
\end{align}
Indeed, the exact nonprincipal source formula is
\[
 \Phi_\chi(s)=
 -\frac{L'(s,\chi)/L(s,\chi)}{s(s-1)}
 +\frac{L'(1,\chi)/L(1,\chi)}{s-1},
 \qquad \Re s>1,
\]
in
\path{Development/AutomorphicArithmetic/DedekindLFunctions/LandauTheorems/mellinGchiOfNeOne.lean},
lines 116--119. The zeta formula, in the adjacent
\path{mellinGchiModOne.lean}, lines 115--118, is the formula used in
\fid{DET-11}, with the additional terms
$-(s-1)^{-2}+(1+\gamma_E)(s-1)^{-1}$.
The change of variable is \texttt{poleData\_change},
\path{poleData.lean}, lines 114--135.
The source conjugation identity
$G_{\chi^\vee}(y)=\overline{G_\chi(y)}$ is
\path{GchiConj.lean}, lines 38--51. It justifies the finite
half-sum used in \eqref{eq:DET18weightedlog}; it is also the exact
half-sum in \path{poleData.lean}, lines 248--261 and 288--304.
Absolute convergence on $\Re s>1$ follows directly by integrating
the absolute ramp majorant, using $|\chi(n)|\le1$ and
$\sum_n\Lambda(n)n^{-\Re s}<\infty$; the linear and $y\log y$
corrections are integrable there as well. The general
\texttt{PoleData} interface states its Laplace identity on the
smaller half-plane $\Re s>2$, and agrees with this formula there.
Finally,
$F_\chi'/F_\chi=a_\chi+L_\chi'/L_\chi$ and
$F_\zeta'/F_\zeta=a_\zeta+\zeta'/\zeta$ prove
\eqref{eq:DET18actualPhi}.

At a completed zero $\rho$, the completion corrections are analytic
and the residue of $\Phi_{\mathcal L}$ is
\[
 -\frac{S_\rho}{\rho(\rho-1)}
 =\frac{S_\rho}{\rho(1-\rho)}.
\]
This recovers exactly the source's coefficient in
\path{poleData.lean}, lines 344--364, including a possible central
zero. It identifies the scalar endpoint after its completion and
polar corrections; it does not identify an upstream period matrix.

\paragraph{When the integer multiplication can be omitted.}
The following is an exact condition, not an assertion that any
actual retained zero has a fractional weight.
\begin{proposition}
There exists $T\in\mathcal S_1$ with ordinary determinant
$D(v)=\det(I-vT)$ such that
\[
 \mathscr L_{\mathcal L}(s)
 =\frac{d}{ds}\log D(s(1-s))
\]
as meromorphic functions if and only if
\begin{equation}\label{eq:DET18integrality}
 S_\rho\in\mathbb Z_{\ge0}\quad(\rho\ne1/2),\qquad
 S_{1/2}/2\in\mathbb Z_{\ge0}.
\end{equation}
If these conditions hold, $T$ can be chosen diagonal with the same
eigenvalues as above and these reduced multiplicities. It is positive
under the stipulated GRH.
\end{proposition}
\begin{proof}
Every zero of an ordinary entire determinant has an integer
multiplicity. At a noncentral preimage the map $s\mapsto s(1-s)$
preserves it, whereas at $s=1/2$ it doubles it. Logarithmic
derivatives have precisely these zero multiplicities as residues.
This proves necessity.

For sufficiency, divide every multiplicity $q_a$ in the product for
$\mathscr D_M$ by $M$. Conditions
\eqref{eq:DET18integrality} say exactly that the results are
nonnegative integers. The reciprocal-eigenvalue sum remains finite,
so the resulting diagonal operator $T$ is trace class. Its
normalized determinant $D$ satisfies $D(v)^M=\mathscr D_M(v)$ by
the absolutely convergent products. Pull back along $v=s(1-s)$
and logarithmically differentiate to obtain the asserted identity.
\end{proof}

Equivalently, the normalized analytic germ
$\mathscr D_M(v)^{1/M}$ at $v=0$ extends to an entire ordinary
determinant precisely when \eqref{eq:DET18integrality} holds.
Positive rational weights alone do not establish these
divisor-integrality conditions. For the actual level, denominators
of $S_\rho$ and $S_{1/2}/2$ divide $2d_{\mathcal L}$. Thus the least
integer multiplication needed at the divisor level is
\[
 \operatorname{lcm}\left(
 \{\operatorname{den}(S_\rho):\rho\ne1/2\}
 \cup\{\operatorname{den}(S_{1/2}/2)\}\right),
\]
where $\operatorname{den}(0)=1$. This least common multiple is
finite because all its entries divide the single integer
$2d_{\mathcal L}$. Its numerical value is not inferred from
positivity.

There is additional structure in the actual source array.
\texttt{KSet} is inverse-closed and its exceptional character is
real, by \path{PrimePairSizeCondition.lean}, lines 39--45.
A real Dirichlet character equals its inverse. Consequently
$I$ is inverse-closed and $\alpha_{\chi^\vee}=\alpha_\chi$:
the exceptional character is fixed by inversion, and all other
characters have the same weight $d_{\mathcal L}^{-1}$.
Reindexing the finite sum then gives
\[
 S_\rho=m_\zeta(\rho)+\sum_{i\in I}\alpha_i m_{\chi_i}(\rho).
\]
In particular, the half in the general family formula is not an
additional denominator at noncentral zeros for this actual array.
This observation still does not evaluate the multiplicities needed
to omit all integer multiplication.

The factor $1/M$ in \eqref{eq:DET18resolvent} also cannot in general
be absorbed by replacing $C_{\mathcal L,M}$ by
$C_{\mathcal L,M}/M$. Under the stipulated GRH, the nonzero positive
operator has $\Tr(C_{\mathcal L,M}^2)>0$. In the expansions at $v=0$,
the coefficients of $v$ in
\[
 \frac1M\Tr(C_{\mathcal L,M}(I-vC_{\mathcal L,M})^{-1})
 \quad\hbox{and}\quad
 \Tr\bigl((C_{\mathcal L,M}/M)
 (I-vC_{\mathcal L,M}/M)^{-1}\bigr)
\]
are respectively $\Tr(C_{\mathcal L,M}^2)/M$ and
$\Tr(C_{\mathcal L,M}^2)/M^2$; they differ when $M>1$.
Integer multiplicities, trace scaling, and eigenvalue scaling are
therefore distinct operations.

This closes the ordinary determinant representation for the
denominator-cleared finite character endpoint, with explicit source
normalizations, full zero multiplicities and central behavior.
The operator is constructed from the zero divisor. The result
supplies neither the historical step-4 operator nor a zero-independent
identification of the earlier finite arithmetic matrices with this
operator.
